%%%%%%%%%%%%%%%%%%%%%%%%%%%%%%%%%%%%%%%%%%%%%%%%%%%%%%%%%%%%
%
% Authors:       Nina Wiedemann, Benjamin Ummenhofer, Diana Wofk, 
%                Quentin Leboutet, and Michael Paulitsch
%
% Description:   The SYCL of Life: Evolutionary Kernel Optimization 
%                with Large Language Models
%
% Section:       Conclusions
%
%%%%%%%%%%%%%%%%%%%%%%%%%%%%%%%%%%%%%%%%%%%%%%%%%%%%%%%%%%%%

\section{Conclusion} \label{sec:conclusion}

We presented KernelFoundry, an evolutionary framework for GPU kernel optimization with LLMs. 
We combine kernel-specific quality-diversity search, meta-prompting and parameter optimization to generate high-performing CUDA and SYCL kernels, demonstrating cross-platform applicability and outperforming SOTA methods.
Our experiments highlight the need for open, reproducible benchmarks and kernel databases to advance the field.\\
Future work includes templated kernel generation for speedups across input ranges and tensor shapes, improved formal verification to eliminate reward hacking, and deeper hardware specialization. Despite the success of prompting-based frameworks, it remains a promising avenue  to finetune models, in particular through RL with verifiable rewards. As LLMs and hardware evolve, frameworks like KernelFoundry will be essential for automated, robust kernel optimization -- enabling rapid deployment of new language models and fast adaptation of kernels to new hardware platforms.
